\section{Theorem 1: multiplication by zero gives zero}\label{sec:10-ring-theorems:02-theorem-1:1}

\textbf{Claim.} \texttt{Rg.mul a Rg.addGrp.id = Rg.addGrp.id}, i.e.~\(a \cdot 0 = 0\).

\textbf{Finding the proof.} In ordinary notation, the standard trick is as follows. It is worth memorizing as a pattern applicable to any additive-identity argument, start from \(0 = 0 + 0\), multiply through, and cancel.

\[
a \cdot 0 \overset{?}{=} 0
\]

No \texttt{Ring} axiom directly states this. \href{https://loogle.lean-lang.org/?q=mul_zero}{\texttt{mul\_\allowbreak{}zero}} is not a field of \texttt{Ring} (Chapter 9); it must instead be \emph{derived} from \texttt{left\_\allowbreak{}distrib} plus group cancellation. The general recipe is that whenever a goal involves \texttt{0} (or any identity element) in a non-trivial way, attempt rewriting it as \texttt{0 + 0}. This is exactly the ``pad with the identity'' trick of \hyperref[sec:08-group-theorems:03-theorem-2:1]{Theorem 2 of Chapter 8}, but here it is applied to the input of the equation to be proved, rather than to the output.

\[
a \cdot 0 = a \cdot (0 + 0) = a\cdot 0 + a \cdot 0
\]

using \texttt{left\_\allowbreak{}distrib} on the last step. Writing \(x := a \cdot 0\), this shows \(x = x + x\). The goal has thereby become ``an additive-group element that equals its own double is zero,'' a pure group fact, provable by adding \(-x\) to both sides and using associativity/cancellation exactly as in Chapter 8. \textbf{Recognizing that the ring-shaped goal reduces to a group-shaped goal already known to be solvable} is the real insight; everything thereafter is mechanical \texttt{rw}.

\begin{lstlisting}
theorem mul_zero (a : R) : Rg.mul a Rg.addGrp.id = Rg.addGrp.id := by
  have h0 : Rg.addGrp.op Rg.addGrp.id Rg.addGrp.id = Rg.addGrp.id :=
    Rg.addGrp.toGroup.id_left Rg.addGrp.id
  have h1 : Rg.mul a (Rg.addGrp.op Rg.addGrp.id Rg.addGrp.id) =
      Rg.addGrp.op (Rg.mul a Rg.addGrp.id) (Rg.mul a Rg.addGrp.id) :=
    Rg.left_distrib a Rg.addGrp.id Rg.addGrp.id
  rw [h0] at h1
  -- h1 : Rg.mul a Rg.addGrp.id = op (mul a 0) (mul a 0), i.e. x = x + x
  have h2 :
      Rg.addGrp.op (Rg.addGrp.toGroup.inv (Rg.mul a Rg.addGrp.id)) (Rg.mul a Rg.addGrp.id) =
      Rg.addGrp.op (Rg.addGrp.toGroup.inv (Rg.mul a Rg.addGrp.id))
        (Rg.addGrp.op (Rg.mul a Rg.addGrp.id) (Rg.mul a Rg.addGrp.id)) :=
    congrArg (Rg.addGrp.op (Rg.addGrp.toGroup.inv (Rg.mul a Rg.addGrp.id))) h1
  rw [Rg.addGrp.toGroup.inv_left] at h2
  rw [(*$\leftarrow$*) Rg.addGrp.toGroup.assoc] at h2
  rw [Rg.addGrp.toGroup.inv_left] at h2
  rw [Rg.addGrp.toGroup.id_left] at h2
  exact h2.symm
\end{lstlisting}

\texttt{h2} is proved with \texttt{congrArg}, not \texttt{by rw [h1]}. In an earlier draft, attempting to rewrite with \texttt{h1} at this intermediate \texttt{have} using plain \texttt{rw} caused occurrence-targeting problems. \texttt{rw [h1]} rewrites \emph{every} syntactic occurrence of the left-hand side of \texttt{h1} in the goal, including copies produced by the substitution itself, and hence does not land on the exact stated goal here. \texttt{congrArg f h1} avoids this problem entirely. It directly builds ``apply \texttt{f} to both sides of \texttt{h1},'' which is precisely ``add \(-x\) to both sides of \(x = x+x\)'' with no ambiguity about which occurrence is targeted.

If progress is lost partway through, the recovery move is always the same, translate the \emph{current} hypothesis (\texttt{h1}, then \texttt{h2}) into ordinary \texttt{+}/\texttt{0} notation, check it against the paper derivation above, and identify exactly which line of the paper proof this corresponds to. The Lean proof is long only because each paper-proof line (``add \(-x\) to both sides'', ``regroup'', ``cancel'') corresponds to one \texttt{rw}. Nothing here is conceptually harder than the five-line sketch preceding it.

\begin{mathreading}

This is the absorbing law \(a\cdot 0 = 0\), which is \emph{not} an axiom but a consequence of distributivity plus additive cancellation. Writing \(x = a\cdot 0\): \[
x = a\cdot 0 = a\cdot(0+0) = a\cdot 0 + a\cdot 0 = x + x,
\] using \(0 = 0+0\) and \texttt{left\_\allowbreak{}distrib}; adding \(-x\) and cancelling gives \(0 = x\). Conceptually this says the map \(x \mapsto a\cdot x\) is a group homomorphism of \((R,+)\), and homomorphisms send the identity to the identity. \(0\) absorbs because multiplication is additive in each argument.

\end{mathreading}

\begin{progcorner}

\texttt{mul\_\allowbreak{}zero} looks obvious enough that a Python codebase would never think to test it, \texttt{x * 0} is \texttt{0}, of course. It genuinely is, for \texttt{int}. It is not, for every numeric type Python ships.

\end{progcorner}

\begin{lstlisting}[language=Python,style=python]
float('nan') * 0.0   # nan, not 0.0
float('inf') * 0.0   # nan, not 0.0
\end{lstlisting}

\texttt{nan} and \texttt{inf} are ordinary \texttt{float} values, reachable from perfectly normal-looking arithmetic (\texttt{1.0 /\allowbreak{} 0.0} under the right settings, or the result of an earlier overflow), and once one of them appears, ``multiply anything by zero and get zero'' quietly stops holding, with no exception raised anywhere to flag it. The theorem \texttt{mul\_\allowbreak{}zero} proved above carries no such asterisk, because it is proved from the \texttt{Ring} axioms of Chapter 9 alone, and any \texttt{R} that can be given a genuine \texttt{Ring R} value must, by that same proof, actually satisfy it, no matter what \texttt{R} turns out to be later. There is no possibility of an \texttt{R} that type-checks as a \texttt{Ring} and still has some element behaving like \texttt{nan}, since a \texttt{nan}- like element would falsify \texttt{left\_\allowbreak{}distrib} or one of the additive-group axioms \texttt{mul\_\allowbreak{}zero} is built from, and a \texttt{Ring R} value could never have been constructed for it in the first place.

\textbf{Mathlib equivalent.} Where the book spends a full \texttt{have}/\texttt{congrArg}/\texttt{rw} derivation getting from \texttt{left\_\allowbreak{}distrib} and group cancellation to \(a\cdot 0=0\), Mathlib already proves this and gives it exactly the same name.

\begin{lstlisting}
example {R : Type*} [Ring R] (a : R) : a * 0 = 0 := mul_zero a
\end{lstlisting}

There is nothing to derive. \texttt{mul\_\allowbreak{}zero} is proved once, generically, for every \texttt{Ring} (indeed every \texttt{MulZeroClass}), by essentially the argument above, and is then simply \emph{available} rather than re-derived at each use site.
